\documentclass[11pt]{article}

\usepackage{amsmath,amsfonts,amssymb}
\usepackage{natbib}
\usepackage{bookstabs}

\title{2. Iris Dataset Set Analysis Using a Voronoi Classifier}

\begin{document}
The Iris data set is a widely distributed collection of data pertaining to three common varieties of Iris,\\
I. setosa, I. versicolor, and I. virginica. For each subspecies, four features where measured:\\
sepal length, sepal width, petal length and petal width.\\

Results: First, a linear classification of the setosa vs versicolor was implimented. The sepal length \\
is plot against sepal width, petal length, and petal width. Then the sepal length vs sepal width is used to \\
classify using both the Gaussian and Sinc kernels. The following is a table of errors determined for each plot.\\
\begin{table}[ht]
\caption{Classification Error on Test Data for each Plot (I. setosa vs I. versicolor)} \\
\begin{tabular}{cc}
\hline
Plot No. & Percent Error \\[0.5ex]
\hline
1  & 0 \\
2  & 5 \\
3  & 0 \\
4  & 0 \\
5  & 10 \\
\end{tabular}
\end{table}
\end{document}

